\section{Related work}

\subsection{Prediction markets and direct probability elicitation}
Prediction markets provide one mechanism for aggregating distributed beliefs, with early organizational applications including sales forecasting and internal corporate markets \citep{chenplott2002,cowgill2009,cowgill2015}. Field evidence also documents departures from an idealized independent-information market, including optimism and correlated trading within organizational and social structures \citep{cowgill2009}. Direct probability elicitation is therefore a substantive comparator rather than a procedural placeholder: prediction polls can retain considerable information and in some settings compete closely with market forecasts \citep{atanasov2017,dana2019}. Study~0 starts with private elicitation because its first question concerns the existence of a distributed engineering signal, not the residual value of a market mechanism.

\subsection{Shared information, meta-predictions and aggregation}
Collective forecasts can be distorted by shared information and correlated errors. A simple probability average is transparent but can be conservative when individual forecasters each condition on only part of the available evidence \citep{baron2014,satopaa2014}. Meta-belief approaches ask respondents not only for their own judgment but also for a prediction of others' judgments. Bayesian Truth Serum introduced second-order population predictions as part of an incentive-compatible elicitation mechanism for subjective data \citep{prelec2004}. The Surprisingly Popular principle later used predicted population responses to recover minority knowledge in single-question crowd-wisdom problems \citep{prelec2017}. Related work uses meta-predictions to adjust for shared information or identify latent expertise without a long individual performance history \citep{palley2019,martinie2020,wilkening2022,peker2025}.

These methods differ in assumptions, targets and algorithms. The fixed logit recalibration in Study~0 is therefore treated as a project-specific, prospectively frozen aggregator inspired by this literature, not as an implementation of Bayesian Truth Serum, Surprisingly Popular, or any cited meta-belief algorithm.

\subsection{Expert judgment in software engineering}
Expert judgment is also an established object of empirical software-engineering research. J{\o}rgensen's review documents extensive use and study of expert estimation for software development effort \citep{jorgensen2004}. That literature motivates taking the accountable engineering owner's judgment seriously as an institutional comparator rather than benchmarking a crowd only against an empty or purely statistical baseline. The target here is different---short-horizon engineering qualification outcomes rather than development effort---so no performance result is transferred from effort-estimation studies.

\subsection{Measurement reactivity and the interference problem}
Elicitation itself need not be behaviorally inert. A meta-analysis of intention and self-prediction questions reports a small average question--behavior effect across behavioral studies \citep{wood2016}. The domains and interventions differ materially from engineering work, so the reported effect size is not imported into this protocol. Instead, the literature supplies a general validity warning: asking an actor to predict an outcome may change attention or behavior. This motivates private elicitation, pre-resolution aggregate suppression, explicit actor/observer metadata and a separate detected-interference sensitivity analysis.

\subsection{Position of the present protocol}
The protocol combines these strands around a narrower engineering question: whether private distributed forecasts add probabilistic information beyond a same-time accountable-owner signal under prospectively resolvable evidence contracts. It does not treat prediction markets, meta-predictions or expert judgment as inherently superior. Its distinguishing object is the coupling of forecast semantics, institutional comparators, authoritative outcome evidence and non-interference safeguards in a single prospective engineering protocol.
